\section{Variant Projection Without a Graph}
\label{sec:deconstruct}

The analyses of \S\ref{subsec:eng-instances} are folds in an obvious sense: they
visit nodes and accumulate. This section develops the one that is not obvious.

Deriving a reference-relative variant table is the most consequential thing done
with a pangenome, because it is the bridge to everything downstream. A pangenome
encodes variation only implicitly, as the alternative routes haplotypes take
through the graph's bubbles, and essentially none of the established genomics
toolchain can read that. Variant annotation, association testing, imputation, and
clinical reporting all consume the Variant Call Format (VCF), so a pangenome
becomes usable to the field at large only after it has been projected into one.

It is also the analysis least suited to streaming, which is why we treat it in
depth. The other five ask questions about node visits. This one asks a question
about \emph{topology}: at which positions does the graph branch, and what does
each haplotype spell as it crosses each branch? The standard tool,
\cmd{vg deconstruct}~\cite{garrison2018variation}, answers it by loading the whole
graph, computing a snarl (superbubble)
decomposition~\cite{onodera2013superbubbles,paten2018superbubbles}, and then, for
each snarl, enumerating the traversals its embedded paths take and emitting those
that differ from the reference. Both the materialized graph and the snarl index
stay resident for the entire run, which is what makes the projection cost
$397$\,GB on an $80$\,GB graph (Section~\ref{sec:eval}).

\subsection{Separating Discovery from Observation}
\label{subsec:dec-idea}

Our design rests on one decomposition: \emph{site discovery and allele observation
need not happen together}.

Where the graph branches is a function of topology alone. It does not depend on
which haplotypes exist, how many there are, or what they spell. What each
haplotype spells at a site, on the other hand, depends only on that haplotype and
on the site's boundaries, and not on any other haplotype. The standard formulation
interleaves the two, which forces both the graph and the decomposition to be
resident at once. Split them, and each half falls into a form the container
already supports.

\begin{itemize}[noitemsep,topsep=3pt,leftmargin=*]
  \item \textbf{Phase 1} reads only the topology columns and runs in $O(V{+}E)$,
        decoding no haplotype. Its output is a list of site boundaries, which is
        small.
  \item \textbf{Phase 2} is an ordinary compressed-traversal fold: each haplotype
        streams past exactly once, through a state machine that carries a few
        words of state.
\end{itemize}

Neither phase materializes the graph, and neither builds a general-purpose
whole-graph snarl index. Before either runs, each reference contig is decoded once
into a node stream, and a prefix sum over the segment length column maps reference
node positions to genomic coordinates.

\subsection{Phase 1: Topology-Only Site Discovery}
\label{subsec:dec-phase1}

\sys defines sites as top-level snarls, matching the record granularity of
\cmd{vg deconstruct}, and recovers them from a biconnected-components
decomposition of the stored edge columns.

The decomposition runs over a \emph{node-end} graph rather than the segment graph.
Each segment $s$ contributes two vertices, its $5'$ and $3'$ ends, joined by a
\emph{black} edge; each GFA link contributes a \emph{grey} edge between the source
departure end and the target arrival end. This is the bidirected interval graph of
snarl theory, in which snarls correspond to biconnected
blocks~\cite{paten2018superbubbles}. The black edges are not bookkeeping: they are
what makes inversions work. A single-node inversion closes through that segment's
two ends and is recovered as one biconnected block, whereas on a collapsed segment
graph the inverted segment would be an articulation point and the site would split
in two.

\sys runs an iterative Hopcroft--Tarjan decomposition over this
graph~\cite{hopcroft1973algorithm}, maintaining discovery times, low-link values,
and an edge stack; on return from child $w$ to parent $u$, the condition
$\mathit{low}(w)\ge\mathit{disc}(u)$ pops one maximal biconnected block. The pass
is iterative rather than recursive because chromosome-scale graphs overflow the
call stack, and it runs in $O(V{+}E)$ time and linear auxiliary space.

Each non-trivial block is then projected onto the reference node stream. If the
first and last reference positions the block touches define a span whose positions
all belong to the block and which repeats no segment, that span is emitted as a
candidate top-level snarl. The filter rejects cyclic and ambiguous reference
traversals, while nested child bubbles and chains of leaf bubbles collapse into
their enclosing top-level snarl, which is the granularity the reference tool
reports. A final reduction leaves a non-overlapping chain ordered by position.

\subsection{Phase 2: Streaming Allele Observation}
\label{subsec:dec-phase2}

With the sites fixed, the projection becomes a fold. Each haplotype streams
through \textsc{StreamNodes} into a state machine that is either closed or
associated with exactly one open site (Algorithm~\ref{alg:snarl}). A boundary
node-side opens a site; ordinary nodes accumulate into a small interior buffer;
the matching exit emits that interior, tagged with its source traversal. An exit
may immediately open the next site, which is what handles chained snarls without
retaining any of the surrounding path.

\begin{algorithm}[t]
\small
\caption{Allele observation for one haplotype}
\label{alg:snarl}
\begin{algorithmic}[1]
  \REQUIRE slice $E$; boundary maps \textsc{Entrance} (node-side $\to$ site,
           forward or reverse) and per-site exits
  \STATE $\textit{open}\gets\bot$;\quad $\textit{interior}\gets[\,]$
  \STATE \textbf{for each} node $u$ decoded from $E$ (Alg.~\ref{alg:stream}):
  \STATE \quad $v\gets$ node-side vertex of $u$
  \STATE \quad \textbf{if} $\textit{open}=\bot$: \textsc{TryOpen}($v$)
         \COMMENT{sets \textit{open}, \textit{inverted}, \textit{exit}}
  \STATE \quad \textbf{else if} $v=\textit{exit}$:
  \STATE \qquad emit \textit{interior} for site \textit{open}
         (reverse-complemented if \textit{inverted}), tagged with $E$
  \STATE \qquad $\textit{open}\gets\bot$; \textsc{TryOpen}($v$)
         \COMMENT{exit may open a chained site}
  \STATE \quad \textbf{else}: append $u$ to \textit{interior}
\end{algorithmic}
\end{algorithm}

Observation runs in parallel over haplotypes, each thread appending tagged
interiors to a private buffer. The buffers are concatenated once, and the
inherently serial per-site assembly and VCF formatting then run as a single pass,
with each site's observations sorted by source slice so that the emitted VCF is
invariant to thread count.

\paragraph{Inversions.}
A haplotype need not cross a site on the reference strand. When it crosses on the
opposite strand, both the boundary node-sides and the interior nodes appear
reverse-complemented, and yet the engine must recognize the same site and recover
the same allele. Phase~1 has already done its part through the black edges. Phase~2
does the rest by giving every site a second, reverse boundary pair: besides the
forward entrance at the reference source exiting at the sink, \sys stores a
reverse entrance at the reverse complement of the sink, exiting at the reverse
complement of the source. A haplotype entering through the reverse entrance is
marked inverted, and at its exit the interior is reverse-complemented before the
allele is spelled, normalizing the observation into the reference frame so that a
forward and an inverted traversal of the same sequence collapse to a single
allele. Note that this normalization is again the compressed representation's own
symmetry (P5) surfacing at the query level: reverse complement is expressible on
the buffer as it is on the rulebook.

\paragraph{VCF conventions.}
Per site, the reference interior spells allele $0$; each observed interior is
spelled by concatenating segment DNA, with a reverse-orientation node
contributing the reverse complement of its sequence, and deduplicated to an allele
index. Sites with no alternate allele are dropped. Equal-length alleles emit as
substitutions; otherwise the record is left-anchored with the last base of the
source anchor, following the standard indel convention. \texttt{INFO} carries
allele counts and frequencies, and an optional guard collapses very large spans to
a symbolic record. Haplotypes collapse into VCF sample columns by the chosen
grouping mode, with one genotype slot per haplotype index, so reference-tiling
contigs resolve to a sample's true ploidy rather than inflating it. Records emit
in sorted genomic order per contig.

\paragraph{Correctness.}
The algorithm is new, but its output is meant to be a drop-in replacement.
Against \cmd{vg deconstruct} on the HPRC chromosome graphs, \sys agrees at
approximately $99.99\%$ of reference positions, producing the same sites and the
same alternate-allele sets except for a handful of records in tangled
pericentromeric regions where the reference traversal itself is ambiguous. We
quantify record-level agreement, alongside time and memory, in
Section~\ref{subsec:eval-deconstruct}.
